% GENERATED FILE. Edit canonical lessons, then run npm run generate:books.
% canonical-source-hash: fnv1a64:c3cbb88374826288
% canonical-lessons: KA-C114-uddesa, KA-C114-yojane, KA-C114-agatya, KA-C114-nirdhara, KA-C114-bayasu

\chapter{Intention, Plan, Need, Decision, To wish for}
\label{ch:ka-intention-plan-need-decision-to-wish-for}
\hlchaptermodality{114}

\begin{chapteropening}
\textbf{By the end of this chapter:} \emph{I can say an intention, a plan, a need, a decision and to wish for.}

Complete the last lesson of chapter 114: I can say an intention, a plan, a need, a decision and to wish for.
\end{chapteropening}

\section[\texorpdfstring{\kn{ಉದ್ದೇಶ} (uddēśa)}{ಉದ್ದೇಶ (uddēśa)}]{\kn{ಉದ್ದೇಶ} (uddēśa) — an intention}

\label{lesson:KA-C114-uddesa}



\begin{quote}
Before the new one: say the Kannada for a wish, then the Kannada for a dream.
\end{quote}

\subsection*{You'll want to know: \kn{ಉದ್ದೇಶ}}
\textbf{\kn{ಉದ್ದೇಶ}} — \emph{uddēśa} — \textquotedblleft{}an intention\textquotedblright{}.

\begin{grammarlens}[title={What you've built}]
One more word for this chapter.
\end{grammarlens}

\subsection*{Your turn}
\begin{itemize}
\raggedright
  \item \emph{Say it:} \emph{uddēśa}
  \item \emph{Say it:} \emph{uddēśa}, once more
  \item \emph{Say it:} say \emph{uddēśa}
  \item \emph{Recall:} say the Kannada for a wish, then the Kannada for a dream, then say \emph{uddēśa} again
\end{itemize}

\begin{tcolorbox}[breakable,colback=teal!4,colframe=teal!35!black,title={Before you move on}]
What does \kn{ಉದ್ದೇಶ} mean? (\textquotedblleft{}An intention\textquotedblright{}.) Say it once more.
\end{tcolorbox}
\section[\texorpdfstring{\kn{ಯೋಜನೆ} (yōjane)}{ಯೋಜನೆ (yōjane)}]{\kn{ಯೋಜನೆ} (yōjane) — a plan}

\label{lesson:KA-C114-yojane}



\begin{quote}
Before the new one: say the Kannada for a dream, then the Kannada for an intention.
\end{quote}

\subsection*{You'll want to know: \kn{ಯೋಜನೆ}}
\textbf{\kn{ಯೋಜನೆ}} — \emph{yōjane} — \textquotedblleft{}a plan\textquotedblright{}.

\begin{grammarlens}[title={What you've built}]
One more word for this chapter.
\end{grammarlens}

\subsection*{Your turn}
\begin{itemize}
\raggedright
  \item \emph{Say it:} \emph{yōjane}
  \item \emph{Say it:} \emph{yōjane}, once more
  \item \emph{Say it:} say \emph{yōjane}
  \item \emph{Recall:} say the Kannada for a dream, then the Kannada for an intention, then say \emph{yōjane} again
\end{itemize}

\begin{tcolorbox}[breakable,colback=teal!4,colframe=teal!35!black,title={Before you move on}]
What does \kn{ಯೋಜನೆ} mean? (\textquotedblleft{}A plan\textquotedblright{}.) Say it once more.
\end{tcolorbox}
\section[\texorpdfstring{\kn{ಅಗತ್ಯ} (agatya)}{ಅಗತ್ಯ (agatya)}]{\kn{ಅಗತ್ಯ} (agatya) — a need}

\label{lesson:KA-C114-agatya}



\begin{quote}
Before the new one: say the Kannada for an intention, then the Kannada for a plan.
\end{quote}

\subsection*{You'll want to know: \kn{ಅಗತ್ಯ}}
\textbf{\kn{ಅಗತ್ಯ}} — \emph{agatya} — \textquotedblleft{}a need\textquotedblright{}.

\begin{grammarlens}[title={What you've built}]
One more word for this chapter.
\end{grammarlens}

\subsection*{Your turn}
\begin{itemize}
\raggedright
  \item \emph{Say it:} \emph{agatya}
  \item \emph{Say it:} \emph{agatya}, once more
  \item \emph{Say it:} say \emph{agatya}
  \item \emph{Recall:} say the Kannada for an intention, then the Kannada for a plan, then say \emph{agatya} again
\end{itemize}

\begin{tcolorbox}[breakable,colback=teal!4,colframe=teal!35!black,title={Before you move on}]
What does \kn{ಅಗತ್ಯ} mean? (\textquotedblleft{}A need\textquotedblright{}.) Say it once more.
\end{tcolorbox}
\section[\texorpdfstring{\kn{ನಿರ್ಧಾರ} (nirdhāra)}{ನಿರ್ಧಾರ (nirdhāra)}]{\kn{ನಿರ್ಧಾರ} (nirdhāra) — a decision}

\label{lesson:KA-C114-nirdhara}



\begin{quote}
Before the new one: say the Kannada for a plan, then the Kannada for a need.
\end{quote}

\subsection*{You'll want to know: \kn{ನಿರ್ಧಾರ}}
\textbf{\kn{ನಿರ್ಧಾರ}} — \emph{nirdhāra} — \textquotedblleft{}a decision\textquotedblright{}.

\begin{grammarlens}[title={What you've built}]
One more word for this chapter.
\end{grammarlens}

\subsection*{Your turn}
\begin{itemize}
\raggedright
  \item \emph{Say it:} \emph{nirdhāra}
  \item \emph{Say it:} \emph{nirdhāra}, once more
  \item \emph{Say it:} say \emph{nirdhāra}
  \item \emph{Recall:} say the Kannada for a plan, then the Kannada for a need, then say \emph{nirdhāra} again
\end{itemize}

\begin{tcolorbox}[breakable,colback=teal!4,colframe=teal!35!black,title={Before you move on}]
What does \kn{ನಿರ್ಧಾರ} mean? (\textquotedblleft{}A decision\textquotedblright{}.) Say it once more.
\end{tcolorbox}
\section[\texorpdfstring{\kn{ಬಯಸು} (bayasu)}{ಬಯಸು (bayasu)}]{\kn{ಬಯಸು} (bayasu) — to wish for, to want}

\label{lesson:KA-C114-bayasu}



\begin{quote}
Before the new one: say the Kannada for a need, then the Kannada for a decision.
\end{quote}

\subsection*{You'll want to know: \kn{ಬಯಸು}}
\textbf{\kn{ಬಯಸು}} — \emph{bayasu} — \textquotedblleft{}to wish for, to want\textquotedblright{}.

\begin{grammarlens}[title={What you've built}]
One more word for this chapter.
\end{grammarlens}

\subsection*{Your turn}
\begin{itemize}
\raggedright
  \item \emph{Say it:} \emph{bayasu}
  \item \emph{Say it:} \emph{bayasu}, once more
  \item \emph{Say it:} say \emph{bayasu}
  \item \emph{Recall:} say the Kannada for a need, then the Kannada for a decision, then say \emph{bayasu} again
\end{itemize}

\begin{tcolorbox}[breakable,colback=teal!4,colframe=teal!35!black,title={Before you move on}]
What does \kn{ಬಯಸು} mean? (\textquotedblleft{}To wish for, to want\textquotedblright{}.) Say it once more.
\end{tcolorbox}
